\begin{corollary}[holomorphic dependence of the comparisons]\label{push:src:E.1.3}
For fixed discrete data in \ref{push:src:E.1.2} and a fixed admissible shift $h$,
the constructed comparisons are holomorphic isomorphisms of families
over a sufficiently small parameter disc.  In particular, for sphere
data with multiplicities $(3,4)$ the identification with the standard
family uses the parameter
\[
 W=e^{-\pi i(c+s_2)}=(-1)^{s_2}w.
\]
The smooth monodromy conjugacy class around the positive parameter
circle is unchanged by this identification of twist data.
\end{corollary}
\Uses{the hypotheses, notation and definitions stated here; \cref{push:src:E.1.2}; \cref{fam:src:B.4.1}}

\begin{proof}
In \ref{push:src:E.1.2} the rational vectors $b_j$, integral defects $\ell_j$,
and toric exponents are independent of $c$.  All period matrices,
root comparisons, and cusp unit factors are holomorphic in $c$.
On a small Stein parameter disc $D$, the correction sheaf is
$\mathcal V_B\boxtimes\mathcal O_D$.  Its first cohomology on
$B\times D$ is zero: project to $D$, use the same two-chart
complex and $H^1(B,\mathcal V_B)=0$, and then Stein acyclicity.
Thus the affine correction can be chosen holomorphically in $c$.
For $(k,l)=(1,1)$ the only nonvacuous finite congruence is
$h\equiv s_2-s_2'\pmod2$, because the admissible $r_2'$ is odd.
Taking $s_2'=0$ gives the displayed parameter identification.
Two choices of these family comparisons differ fibrewise by an
automorphism of the standard member.  The full automorphism
calculation re-derived in the credited appendix at \ref{fam:src:B.4.1} identifies that group with $\mathbb C^*$;
it is connected.  On overlapping parameter arcs the transitions
are therefore isotopic to the identity.  Composing around a circle
changes the resulting diffeomorphism only by such an isotopy and
by conjugation at the chosen base fibre.  Multiplication of the
parameter by $(-1)^{s_2}$ preserves the positive orientation.
\end{proof}